\begin{textsample}{NEG Dim 1 – persona\_gemini – Score -3.00 – t350\_gemini.txt}  \label{ex:f1_neg_001}
Sustainability certification labels found on products like cocoa, coffee, \textbf{palm} oil, soy, and wood often fail to deliver on their promises.
Schemes such as the Forest \textbf{Stewardship} Council ( FSC ), the Programme for the Endorsement of Forest Certification ( PEFC ), and the Rainforest Alliance are not preventing deforestation, the destruction of ecosystems, or human rights abuses.
These \textbf{violations} include infringements on the rights of Indigenous peoples and workers.
Despite their popularity since the 1990s, these labels allow companies to continue their harmful practices while misleading consumers.
The responsibility for change should not be placed on individual buyers.
Instead, systemic change is needed to fight the interconnected crises of climate, biodiversity, and inequality.
Governments in both producer and consumer countries must step up by implementing stronger regulations and effective tracking systems.
Companies, for their part, must adopt robust standards.
The focus must shift from consumer choice to corporate and governmental accountability.

% matched lemmas: palm, stewardship, violation
\end{textsample}
